\section{Spolupráce v týmu}

\begin{table}[h]
    \centering
    \begin{tabular}{p{0.28\textwidth}p{0.62\textwidth}}
        \textbf{Jakub Svoboda}      & sémantický analyzátor, dokumentace, nastavení git repozitáře, testy \\
        \textbf{Štěpán Zobal}       & lexikální analyzátor, dokumentace                                   \\
        \textbf{Matěj Tajovský}     & generátor kódu, dokumentace                                         \\
        \textbf{Sebastián Ondruška} & syntaktický analyzátor, návrh komunikace modulů, dokumentace        \\
    \end{tabular}
\end{table}

Na začátku byla dohoda, že se vytvoří hlavičkové soubory jednotlivých modulů co nejdříve,
aby mohl každý pracovat na své odpovědné části.
Pull requesty byla snaha tvořit vždy, když se implementovala značná funkční část modulu.
Pro pull request bylo nutné, aby byl po code review schválen alespoň jedním členem,
který se na pull requestu nepodílel.

Testování probíhá jako unit testy a integrační testy překladače jako celku.
Unit testy kontrolují jednotlivé funkce každého modulu. Testují ideální použití i krajní případy nebo chybové použití.
Integrační testy probíhají tak, že se přeloží kód, správný nebo chybový, poté se spustí výsledný kód interpretem ic26int a
porovná se výstup s očekávaným výstupem. Pro unit testy se používá testovací framework pro jazyk C cmocka a integrační testy
se provádí pomocí bash skriptu a nástroje diff.

Zvolený styl formátování C kódu je pomocí clang-format verze v18.1.8 skrze nástroj pre-commit.
